\begin{abstract}
This report studies whether \emph{class-conditional synthetic images} generated by a Conditional Generative Adversarial Network (CGAN) can improve downstream image classification when real labeled data are limited. A three-class fruit image dataset (apple, banana, orange) is used as a controlled benchmark. A CGAN with a \textbf{conditional batch-normalized generator} and a \textbf{projection discriminator} is trained to produce 64×64 RGB images. Generator checkpoints are evaluated with the \textbf{Fréchet Inception Distance (FID)} on a held-out validation split, and a synthetic image pool is generated from the best-performing checkpoint. The downstream utility of synthetic data is assessed by training a compact convolutional classifier (FruitCNN) under three scenarios, real-only, synthetic-only, and real+synthetic, across multiple data sizes (100--1300 images per class).

Results show that the CGAN achieves its best recorded validation FID at epoch 140 (\textbf{FID ≈ 143.36}). For classification, real-only training already yields near-ceiling accuracy (≈98.6--99.9\%) even with 100 images per class, reflecting an ``easy'' dataset regime. Adding synthetic images sometimes improves or matches performance (for example, at 400 images/class, real+synthetic reaches \textbf{99.66\%} vs \textbf{98.37\%} for real-only), but provides limited gains when real data are abundant. Training on synthetic data alone generalizes poorly to real test images (accuracy drops to \textbf{69.58\%} at 1300 images/class), indicating a substantial domain gap and underscoring that low FID does not guarantee downstream utility. The report concludes with limitations (single seed, small class set, single architecture) and a roadmap for more robust evaluation (multi-seed experiments, stronger GAN baselines, per-class metrics, and improved synthetic--real alignment).
\end{abstract}

\section*{Executive Summary}
\addcontentsline{toc}{section}{Executive Summary}

\textbf{Goal.} Determine whether CGAN-generated synthetic images can supplement real fruit images to improve classification performance under varying data availability.

\textbf{Approach.}
\begin{itemize}
\item Train a \textbf{conditional GAN} (CGAN) on a 3-class fruit dataset (apple/banana/orange).
\item Monitor training with \textbf{FID} and qualitative sample grids.
\item Generate a balanced synthetic pool from the best-FID checkpoint.
\item Train the same classifier across a \textbf{data-size grid} in three scenarios: real-only, synthetic-only, and real+synthetic.
\end{itemize}

\textbf{Key findings.}
\begin{itemize}
\item \textbf{GAN quality trend:} FID decreases substantially over training and reaches its best value at epoch 140 (FID ≈ 143.36), then degrades slightly by epoch 200.
\item \textbf{Downstream utility:} Synthetic augmentation provides \textbf{limited benefits} because real-only accuracy is already very high at small data sizes on this dataset.
\item \textbf{Domain gap:} Synthetic-only training does not transfer well to real test images, suggesting that the generator distribution differs meaningfully from the real distribution despite improvements in FID.
\end{itemize}

\textbf{Takeaway.} Synthetic data are most credible as a \emph{supplement} to real data, not a replacement, and must be evaluated with downstream tasks and robustness checks (multi-seed, stronger baselines) rather than relying on a single image-quality metric.

\textbf{Reproducibility.} All figures and tables used in this report are exported into \code{reports/figures/} and \code{reports/tables/} and can be regenerated from repository artifacts (Appendix~\ref{app:reproducibility}).

\clearpage
